\subsubsection{Inner Product Space}

\textit{Note:The development of concepts need to keep extension to infinite analog in mind}

\paragraph{The identity operator and the Kroncker's Delta}

Let $V$ be a vector space endorsed with an inner product $\langle\cdot| \cdot \rangle$ of dimension $n$ and $\beta = \{e_\nu\}$ be an orthonormal basis. Then we can compute the inner product between arbitrary $x,y\in V$ as follows:
\begin{equation}
\langle x | y \rangle  &= \langle e_{\nu} | x \rangle^* \langle e_{\nu} | y\rangle \\ 
&= \langle x | e_{\nu} \rangle \langle e_{\nu} | y\rangle \\
&= \langle x \lvert |e_{\nu} \rangle \langle e_{\nu}|  \rvert y\rangle
\end{equation}
But $\langle x | y \rangle = \langle x | I | y \rangle$, hence we get that

\begin{lemma}[Outer-product decomposition of the identity operator]Let $\{e_{\nu}\}$ be an orthonormal basis for $V$, then:
\begin{equation}
|e_{\nu} \rangle \langle e_{\nu}|=I
\end{equation}

\end{lemma}\begin{framed}
\textbf{\textit{Help!} What is happening here?}\\
\textit{Step 1} We first rewrite inner product as a sum of component-component product. In $\mathbb{C}^3$ for instance,
\begin{equation}
\langle x | y\rangle &= \sum_{\nu} (x^{\nu})^* y^{\nu} \\
&= \sum_{\nu} \langle e_{\nu} | x \rangle^*  \langle e_{\nu} | y \rangle
\end{equation}
Note in the last step we envicted orthonormality of $\{e_\nu\}$. Adopting Einstein notaiton we get the first equality.

\textit{Step 2} We are bascially writing two consecutive inner product $\langle x, \alpha \rangle \langle \beta, y \rangle$ as a bilinear form $\langle x, A y\rangle$ where $A=| \alpha \rangle \langle \beta |$. Consider this example in $\mathbb{R^3}$:

\begin{equation}
\langle x, \alpha \rangle \langle \beta, y \rangle &=
\begin{pmatrix}
x^1 & x^2 & x^3
\end{pmatrix}
\begin{pmatrix}
\alpha^1 \\ \alpha^2 \\ \alpha^3
\end{pmatrix}
\begin{pmatrix}
\beta^1 & \beta^2 & \beta^3
\end{pmatrix}
\begin{pmatrix}
y^1 \\ y^2 \\ y^3
\end{pmatrix} \\
&= 
\begin{pmatrix}
x^1 & x^2 & x^3
\end{pmatrix}
\begin{pmatrix}
\alpha^1 \beta^1 & \cdots & \alpha^1\beta^3 \\
\vdots & & \vdots \\
\alpha^3 \beta^1 & \cdots & \alpha^3\beta^3
\end{pmatrix}
\begin{pmatrix}
y^1 \\ y^2 \\ y^3
\end{pmatrix} \\
\end{equation}

In otherword, we are just invoking associative law(which holds in vector space):
\begin{equation}
(x\alpha^T) (\beta y^T)&=x(\alpha^T \beta) y^T
\end{equation}
where the outer product $\alpha^T \beta=A$.

We denote bilinear forms $x^TAy$ in sandwich notations: $\langle x |A|y\rangle$

\textit{(Step 3)} The only possble opeartor to be sandwitched without changing the result:
\begin{equation}
\langle x|Ay\rangle &= \langle x | y\rangle
\end{equation}
is when $A=I$. We therefore conclude that the outer product $|e_{\nu}\rangle \langle e_{\nu}|$ is the identity. In math notation, this is saying:

\begin{itemize}
\item For orthonormal basis $\{e_{\nu}\}$, we have

\begin{equation}
\sum_{\nu=1}^n e_{\nu}e_{\nu}^T=I
\end{equation}
\end{itemize}
\end{framed}

Now we consider representing $|x\rangle$ in $\beta$:
\begin{equation}
|x\rangle &= (|e_{\nu} \rangle \langle e_{\nu}|) |x\rangle \\
&= |e_{\nu}\rangle \langle e_{\nu}|x\rangle
\end{equation}
Now consider taking the inner product with another basis $e_{\mu}$:
\begin{equation}
\langle e_{\mu}|x\rangle  &=\langle e_{\mu}|e_{\nu}\rangle \langle e_{\nu}|x\rangle
\end{equation}